\chapter{Manuscript statements and Lean declarations}
\label{chap:current-statements}

This table gives an entry point for each numbered manuscript statement.
The selected declarations include definitions, deductions and explicit
assumptions. Their association with a statement does not assert that every
part of its proof is formalised or that its external hypotheses have been
established. The relevant chapters explain those hypotheses and the full
types in Lean specify them.

The source associations accompanying the package list the additional
declarations used by each argument. Numbers within declaration names are
not manuscript cross-references.

\begingroup
\small
\setlength{\tabcolsep}{4pt}
\renewcommand{\arraystretch}{1.04}
\begin{longtable}{@{}>{\raggedright\arraybackslash}p{0.13\textwidth}>{\raggedright\arraybackslash}p{0.49\textwidth}>{\raggedright\arraybackslash}p{0.32\textwidth}@{}}
\toprule
Manuscript & Selected declaration & Role and scope \\
\midrule
\endfirsthead
\toprule
Manuscript & Selected declaration & Role and scope \\
\midrule
\endhead
\bottomrule
\endfoot
Theorem\par 1.1 & \lean{ManuscriptIBAW.Main.theorem_1_1} & Conditional deduction for the three families. \\
\addlinespace[2pt]
Corollary\par 1.2 & \lean{ManuscriptIBAW.Main.corollary_1_2} & Deduction using the assumed finite group reduction. \\
\addlinespace[2pt]
Definition\par 2.1 & \lean{ManuscriptIBAW.Preliminaries.CentralSectors.existsUnique_centralCharacterSector} & Existence and uniqueness of the central character sector. \\
\addlinespace[2pt]
Lemma\par 2.2 & \lean{ModularRep.LinearBlockStabilizer.modular_block_stabilizer_bound} & Modular stabiliser bound under the block and character hypotheses. \\
\addlinespace[2pt]
Theorem\par 2.3 & \lean{ManuscriptIBAW.Preliminaries.SpathCriterion.current_theorem_2_3} & Character triple equivalence under an explicit source assumption. \\
\addlinespace[2pt]
Remark\par 2.4 & \lean{ManuscriptIBAW.Preliminaries.BlockFamilies.fullFamily} & Part (a): construction of the family using the assumed passage from block conditions to the full condition. Part (b) is an orbit transport observation. \\
\addlinespace[2pt]
Corollary\par 2.5 & \lean{ManuscriptIBAW.Preliminaries.Descent.current_corollary_2_5_coherent} & Descent with explicit character and root compatibility. \\
\addlinespace[2pt]
Lemma\par 2.6 & \lean{ManuscriptIBAW.Preliminaries.CentralCore.core_le_raw} & Normal core containment in the central quotient construction. \\
\addlinespace[2pt]
Corollary\par 2.7 & \lean{ModularRep.PaperProofs.ConlonBasicSet.corollary_2_4_conlonMark} & Conlon deduction for the stated prime and lattice assumptions. \\
\addlinespace[2pt]
Lemma\par 2.8 & \lean{ManuscriptIBAW.equivariant_bijection_of_basic_set_fixed} & Equivariant bijection from a fixed integral basic set. \\
\addlinespace[2pt]
Lemma\par 2.9 & \lean{ManuscriptIBAW.Jordan.SemidirectExtension.restrict} & Restriction of an extension from a semidirect product. \\
\addlinespace[2pt]
Lemma\par 2.10 & \lean{ManuscriptIBAW.TypeC.EvenApplication.unipotent_bawGood} & Application of the cyclic outer criterion under its source assumptions. \\
\addlinespace[2pt]
Lemma\par 2.11 & \lean{ManuscriptIBAW.FiniteCharacterSelection.exists_delta_injection_of_nonnegative} & Finite character selection. Geometry and trace identities remain inputs. \\
\addlinespace[2pt]
Lemma\par 2.12 & \lean{ManuscriptIBAW.TypeC.GGGRSelection.delta_selection} & Conditional Type C selection from the specified geometric sources. \\
\addlinespace[2pt]
Theorem\par 3.1 & \lean{ManuscriptIBAW.TypeC.all_primes} & Conditional deduction for the Type C cases. \\
\addlinespace[2pt]
Lemma\par 3.2 & \lean{ManuscriptIBAW.Jordan.generalPerLabel_of_fullActor} & Restriction of the Brauer character hypothesis to the automorphism groups chosen for individual semisimple parameters. \\
\addlinespace[2pt]
Proposition\par 3.3 & \lean{ManuscriptIBAW.TypeC.prop33_originalIBAW} & Principal block application with explicit field, counting and compatibility inputs. \\
\addlinespace[2pt]
Remark\par 3.4 & \lean{ModularRep.PaperProofs.OddTwoWreathCoreCharacterSource.triangular_ne_two} & Arithmetic distinction only. The subgroup identifications remain external. \\
\addlinespace[2pt]
Proposition\par 3.5 & \lean{ManuscriptIBAW.TypeC.prop34_full_block_condition} & Conditional rank two conclusion under the sources for the selected prime and field. \\
\addlinespace[2pt]
Lemma\par 3.6 & \lean{ModularRep.PaperProofs.EvenFieldLemmas35_36Actual.lemma_3_5_actual} & Field invariance of the specified ordinary characters and generic weights under the source assumptions. \\
\addlinespace[2pt]
Lemma\par 3.7 & \lean{ModularRep.PaperProofs.EvenFieldLemmas35_36Actual.lemma_3_6_actual} & Equivariant correspondence between the ordinary character set and generic weights under the source assumptions. \\
\addlinespace[2pt]
Proposition\par 3.8 & \lean{ManuscriptIBAW.TypeC.EvenApplication.unipotent_definition35} & Unipotent block condition under the specified published criterion. \\
\addlinespace[2pt]
Proposition\par 3.9 & \lean{ManuscriptIBAW.TypeC.EvenApplicationRouting.fullHG_strictBlocks_of_proposition39_semisimple} & Jordan reduction from the semisimple parameter hypotheses. \\
\addlinespace[2pt]
Lemma\par 3.10 & \lean{ManuscriptIBAW.quotientKernelCenterEquiv} & Quotient and kernel comparison used in the stabiliser bound. \\
\addlinespace[2pt]
Proposition\par 3.11 & \lean{ManuscriptIBAW.TypeC.OddConformalCorrespondence.globalCorrespondence} & Conformal correspondence under the source and action assumptions. \\
\addlinespace[2pt]
Lemma\par 3.12 & \lean{ManuscriptIBAW.TypeC.OddGlobalFactorization.GlobalBijection.allBrauerFactorization} & Transfer of ordinary stabiliser factorisation to Brauer characters. \\
\addlinespace[2pt]
Remark\par 3.13 & \lean{ManuscriptIBAW.TypeC.OddGlobalFactorization.BlockSourceInputs} & Record of the assumptions used for the Conlon argument. It supplies neither canonical labels nor unitriangularity. \\
\addlinespace[2pt]
Proposition\par 3.14 & \lean{ManuscriptIBAW.TypeC.OddPrimeApplication.full_block_condition_source_instantiated} & Application of the full criterion under the listed sources. \\
\addlinespace[2pt]
Theorem\par 4.1 & \lean{ManuscriptIBAW.TypeB.AllPrimes.all_primes} & Conditional deduction for the Type B cases. \\
\addlinespace[2pt]
Lemma\par 4.2 & \lean{ManuscriptIBAW.TypeB.OddConlon.block_bijection} & Blockwise Conlon correspondence under the stated basic set hypotheses. \\
\addlinespace[2pt]
Proposition\par 4.3 & \lean{ManuscriptIBAW.TypeB.OddUniform.complete} & Odd prime application with its structural and criterion sources. \\
\addlinespace[2pt]
Lemma\par 4.4 & \lean{ManuscriptIBAW.TypeB.Levi.component_return_full} & Action on a factor after a cycle of the factor permutation. \\
\addlinespace[2pt]
Lemma\par 4.5 & \lean{ManuscriptIBAW.TypeB.Levi.cartesian_character_orbit} & Character orbit calculation for the rational product. \\
\addlinespace[2pt]
Lemma\par 4.6 & \lean{ManuscriptIBAW.TypeB.Levi.lemma46} & Clifford stabiliser deduction on the specified characters. \\
\addlinespace[2pt]
Lemma\par 4.7 & \lean{ManuscriptIBAW.TypeB.PrincipalField.rationalClass_fixed} & Field invariance of the specified rational classes. \\
\addlinespace[2pt]
Lemma\par 4.8 & \lean{ManuscriptIBAW.TypeB.PrincipalSeries.lemma_4_8} & Principal series deduction under its source interpretation. \\
\addlinespace[2pt]
Proposition\par 4.9 & \lean{ManuscriptIBAW.TypeB.GGGRBasis.gggr_basis} & Conditional basis of projected generalised Gelfand--Graev characters from the geometric and finite selection inputs. \\
\addlinespace[2pt]
Corollary\par 4.10 & \lean{ManuscriptIBAW.TypeB.PrincipalApplication.principal_selector} & Principal character selection using field invariance and extension inputs. \\
\addlinespace[2pt]
Proposition\par 4.11 & \lean{ManuscriptIBAW.TypeB.Exceptional.Inputs.complete} & Exceptional cover application with computation, quotient and root assumptions. \\
\addlinespace[2pt]
Lemma\par 4.12 & \lean{ManuscriptIBAW.Jordan.TypeBLocalized.full_family_of_hypotheses} & Localised Jordan reduction under its factor and transport hypotheses. \\
\addlinespace[2pt]
Proposition\par 4.13 & \lean{ManuscriptIBAW.TypeB.TwoApplication.Inputs.brauerHypothesis} & Derivation of the required Brauer character hypothesis. \\
\addlinespace[2pt]
Proposition\par 4.14 & \lean{ManuscriptIBAW.TypeB.TwoApplication.Inputs.complete} & Prime two application under the principal and Levi sources. \\
\addlinespace[2pt]
Corollary\par 4.15 & \lean{ManuscriptIBAW.TypeB.finite_simple_typeB} & Conditional deduction for the specified Type B groups. \\
\addlinespace[2pt]
Theorem\par 5.1 & \lean{ManuscriptIBAW.Sporadic.Theorem.Inputs.complete} & Sporadic deduction including the external result for 22 groups. \\
\addlinespace[2pt]
Lemma\par 5.2 & \lean{ManuscriptIBAW.Sporadic.CompleteCollapse.lemma_5_2} & Numerical construction with compatible roots for the constructed family. \\
\addlinespace[2pt]
Lemma\par 5.3 & \lean{ManuscriptIBAW.Sporadic.Cancellation.cancel_equivariant_equiv} & Equivariant cancellation of the specified parts. \\
\addlinespace[2pt]
Proposition\par 5.4 & \lean{ManuscriptIBAW.Sporadic.Theorem.BoundaryInputs.proposition_5_4} & Four remaining groups under their numerical and compatible packet sources. \\
\addlinespace[2pt]
\end{longtable}
\endgroup
